\chapter{人身犯罪}

% Source PDF physical page 220; see audit/source_coverage.json.
\begin{EditionOpeningPanel}
1.具体罪名的知识体系：保护法益—犯罪构成要件一法定刑升格条件一犯罪形态一共同犯罪一罪数（罪与罪的区分及联系）。每个罪名遇到的问题不同，所以上述知识体系中的知识点会有详有略。

2.重要考点：非法拘禁罪，绑架罪，拐卖妇女、儿童罪。
\end{EditionOpeningPanel}

\section{侵犯生命、身体的犯罪}

[问题]甲的母亲身患绝症，非常痛苦，让甲找医生给自己打一针送自己上路。甲答应并照办。对甲和医生该如何处理？（安乐死案）

\subsection*{一、故意杀人罪}

第232条故意杀人的，处死刑、无期徒刑或者十年以上有期徒刑；情节较轻的，处三年以上十年以下有期徒刑。

故意杀人罪，是指故意非法剥夺他人生命的行为。本罪侵犯的法益是他人的生命权。

1.安乐死问题

安乐死，是指在患者距离死亡尚有一段时间时，经患者承诺，对其注射无痛苦的毒针，致其死亡。这种行为构成作为的故意杀人罪，因为承诺放弃生命是无效的。例如，本节开头的“安乐死案”中，甲和医生构成故意杀人罪。（2014年第15题）

2.自杀、自伤问题的总结

\InsertDiagram{figures/manual_table_p220_1.tex}

% Source PDF physical page 221; see audit/source_coverage.json.
\ContinuedTable{figures/manual_table_p221_1.tex}

\subsection*{二、故意伤害罪}

第234条故意伤害他人身体的，处三年以下有期徒刑、拘役或者管制。

犯前款罪，致人重伤的，处三年以上十年以下有期徒刑；致人死亡或者以特别残忍手段致人重伤造成严重残疾的，处十年以上有期徒刑、无期徒刑或者死刑。本法另有规定的，依照规定。

故意伤害罪，是指故意非法损害他人身体健康的行为。本罪侵犯的法益是他人的身体健康权。

1.构成要件

（1）行为主体。故意伤害致人轻伤的行为主体是已满16周岁的自然人，故意伤害致人重伤或者死亡的行为主体是已满12或14周岁的自然人。

（2）行为对象是他人的身体。自伤自残不成立本罪，但是战时自伤构成战时自伤罪（第434条）。故意毁坏他人的假肢、假牙、隐形眼镜等，没有造成他人身体伤害的，不构成故意伤害罪，而可能构成故意毁坏财物罪。

（3）伤害行为。不是任何暴行行为都是故意伤害罪中的伤害行为。

我国将伤害结果分为轻微伤、轻伤、重伤与伤害致死。造成轻微伤只需依《治安管理处罚法》处理。造成轻伤结果就由刑法处理。重伤的标准是造成肢体残疾、丧失器官机能（如丧失听觉、视觉）。

伤害行为的成立条件：第一，主观心理具有造成轻伤的意图。第二，客观行为具有造成轻伤的一般可能性。轻伤实际已经是比较严重的伤害，例如，砍掉小手指才算轻伤。例如，甲乙因为自行车剐蹭发生口角，进而推，甲推了乙一把，乙头部倒地，心脏病发作死亡。甲的行为不属于故意伤害罪的伤害行为，如果有过失，可定过失致人死亡罪，但不能定故意伤害罪致人死亡。实务中常见错误是，根据结果来反推行为，如果结果是轻伤、重伤或死亡，就认为行为构成故意伤害罪。

% Source PDF physical page 222; see audit/source_coverage.json.
2.区分：故意杀人罪、故意伤害罪（致人死亡）、过失致人死亡罪。

（1）故意杀人罪与故意伤害罪（致人死亡）的区分。

区分标准：前罪对死亡结果是故意心理，后罪对死亡结果是过失心理。判断方法：看客观证据，例如，打击部位是否致命、是否使用凶器、危害行为是否有节制等。

提示】第一，不能因为行为人没有蓄意谋杀，就认为行为人没有杀人故意。临时起意杀人的情形很多。第二，主观不计后果、不顾死伤的，按照实际结果来定，造成死亡的，定故意杀人罪；造成伤害的，定故意伤害罪。这是因为两种后果都在行为人故意范围内，属于概括故意。（2014年第15题）

（2）故意杀人罪、故意伤害罪（致人死亡）与过失致人死亡罪的区分。

区分标准：过失致人死亡罪无故意杀人行为、故意伤害行为。

3.罪数问题

（1）法条竞合：故意杀人罪与故意伤害罪的包容关系。

故意杀人与故意伤害是性质相同、程度不同的两种行为，二者不是A与-A的对立排斥关系，而是A与A+B的包容评价关系，也即故意杀人可以包容评价为故意伤害。换言之，二者是法条竞合关系，触犯故意杀人罪必然触犯故意伤害罪。

例如（2015年第16题），甲以伤害故意砍乙两刀，随即心生杀意又砍两刀，但四刀中只有一刀砍中乙并致其死亡，且无法查明由前后四刀中的哪一刀造成死亡。杀害行为可以包容评价为伤害行为。因此，可以认为甲前后均实施的是伤害行为，该伤害行为导致死亡结果，因此甲构成故意伤害罪致人死亡。\footnote{从另一个角度也可以分析。具体而言，根据存疑时有利于被告原则，死得越早，对被告人越有利。认定是伤害时的两刀导致死亡，甲构成故意伤害罪致人死亡。}

（2）想象竞合。

常考罪名：暴力危及飞行安全罪，强制猥亵、侮辱罪，非法拘禁罪，侮辱罪，暴力干涉婚姻自由罪，妨害公务罪，抗税罪，强迫卖血罪。在实施这些罪时，同时又故意伤害致人重伤的，想象竞合，择一重罪论处。一般以故意伤害罪（重伤）论处。

\subsection*{三、组织出卖人体器官罪}

第234条之一组织他人出卖人体器官的，处五年以下有期徒刑，并处罚金；情节严重的，处五年以上有期徒刑，并处罚金或者没收财产。

未经本人同意摘取其器官，或者摘取不满十八周岁的人的器官，或者强迫、欺骗他人捐献器官的，依照本法第二百三十四条（故意伤害罪）、第二百三十二条（故意杀人罪）的规定定罪处罚。

违背本人生前意愿摘取其户体器官，或者本人生前未表示同意，违反国家规定，违背其近亲属意愿摘取其户体器官的，依照本法第三百零二条（盗窃、侮辱、故意毁坏户体、尸骨、骨灰罪）的规定定罪处罚。

1．“人体器官”

（1）本罪的“人体器官”主要考虑可移植性，包括眼角膜、皮肤、骨头，不包括血液、骨髓、细胞。（2）本罪的“人体器官”，只包括活体器官，不包括尸体器官。本罪第3款是关于户体的规定。（2011年第14题）

% Source PDF physical page 223; see audit/source_coverage.json.
2.“出卖”

（1）出卖者（器官提供者）不构成本罪。

(2）出卖行为必须是有效的承诺。（2012年第16题，2014年第15题，2016年第58题)

\InsertDiagram{figures/manual_p223_1.tex}

本罪第2款是关于被害人承诺的规定。\footnote{第一，关于承诺限度问题，有人可能问：假如，器官提供者基于真实意愿出卖器官，组织者给提供者造成重伤，由于重伤不能承诺放弃，因此提供者的承诺无效，那么组织者应构成故意伤害罪？对此回答是，原则上，重伤不能承诺放弃，但为了保护另一法益，比如，舍已救人，或将器官移植给另一患者，则重伤可以承诺放弃，在此应当尊重承诺者的自我决定权。第二，本罪的保护法益是合法的器官捐献制度（社会法益）。本罪打击的是黑市上的器官买卖。}

3.“组织”

（1）被组织者，既可以是多人，也可以是一人。也即，组织一人出卖人体器官的，也成立本罪。

（2）组织行为，包括招募、引诱、介绍行为。

（3）单纯购买者、器官接受者，不构成本罪。如果有组织行为，则构成本罪。

4.主观上，不要求行为人具有营利目的。（2011年第14题）

\subsection*{四、遗弃罪}

第261条对于年老、年幼、患病或者其他没有独立生活能力的人，负有扶养义务而拒绝扶养，情节恶劣的，处五年以下有期徒刑、拘役或者管制。

1.行为主体：负有扶养义务的人

本罪的行为主体不限于家庭成员。扶养义务应扩大解释为扶助义务、救助义务。是否有救助义务，可根据总论不作为犯中的作为义务来源来判断。例如，孤儿院、养老院、精神病院的管理人员对所收留的孤儿、老人、精神病人具有扶助义务。某孤儿院院长用卡车将低能儿抛弃到市郊，构成遗弃罪。

2.不作为

本罪是真正不作为犯。不履行扶养义务的方式有：例1，甲将婴儿放到民政局门前。

例2，甲离家出走，不管家中的婴儿。例3，甲开车不慎撞伤乙，驾车逃逸，不管受伤倒地的乙。可见，不履行扶养义务（作为义务）的方式，既有积极的举动，也有消极的静止。但不能因此认为，“遗弃罪既可以由作为构成，也可以由不作为构成，所以是不真正不作为犯”。

3.遗弃罪与不作为的故意杀人罪的关键区分：生命所面临的危险是否紧迫，生命对作为义务的依赖程度，行为人履行义务的难易程度。例如，甲将婴儿置于医院门前，属于遗弃罪。乙将婴儿放在空旷无人的高楼楼顶上，属于不作为的故意杀人罪。（2011年第13题）

% Source PDF physical page 224; see audit/source_coverage.json.
[注意】两罪的区分不是A与-A的对立排斥关系，而是程度之分，也即行为人构成不作为故意杀人罪不意味着就不构成遗弃罪。例如，将婴儿遗弃在荒山野岭，既构成不作为的故意杀人罪，也构成遗弃罪，最终以故意杀人罪论处。

\OriginalPanel{典型真题}

关于故意杀人罪、故意伤害罪的判断，下列哪一选项是正确的？（2014年·卷二·15题）\footnote{[答案]B。}

A.甲的父亲乙身患绝症，痛苦不堪。甲根据乙的请求，给乙注射过量镇定剂致乙死亡。乙的同意是真实的，对甲的行为不应以故意杀人罪论处

B.甲因口角，捅乙数刀，乙死亡。如甲不顾乙的死伤，则应按实际造成的死亡结果认定甲构成故意杀人罪，因为死亡与伤害结果都在甲的犯意之内

C.甲谎称乙的女儿丙需要移植肾脏，让乙捐肾给丙。乙同意，但甲将乙的肾脏摘出后移植给丁。因乙同意捐献肾脏，甲的行为不成立故意伤害罪

D.甲征得乙（17周岁）的同意，将乙的左肾摘出，移植给乙崇拜的歌星。乙的同意有效，甲的行为不成立故意伤害罪

\section{侵犯性权利的犯罪}

[问题]卖淫女同意与甲男、乙男发生性行为，但要求分别进行，而甲乙不答应，强行同时发生性行为。甲乙是否构成强奸罪？（强行玩3P案）

\subsection*{一、强奸罪}

第236条以暴力、胁迫或者其他手段强奸妇女的，处三年以上十年以下有期徒刑。

奸淫不满十四周岁的幼女的，以强奸论，从重处罚。

强奸妇女、奸淫幼女，有下列情形之一的，处十年以上有期徒刑、无期徒刑或者死刑：

\subsubsection*{（一）强奸妇女、奸淫幼女情节恶劣的；}

\subsubsection*{（二）强奸妇女、奸淫幼女多人的；}

\subsubsection*{（三）在公共场所当众强奸妇女、奸淫幼女的；}

\subsubsection*{（四）二人以上轮奸的；}

\subsubsection*{（五）奸淫不满十周岁的幼女或者造成幼女伤害的；}

\subsubsection*{（六）致使被害人重伤、死亡或者造成其他严重后果的。}

强奸罪，是指违背妇女意志，使用暴力、胁迫或者其他手段，强行与妇女发生性交的行为。

\subsubsection*{（一）保护法益}

本罪的保护法益是妇女的性自主权。性自主权包括：（1）是否发生性交的决定权；

（2）关于性交具体条件的决定权，例如，性交对象、时间、地点等条件的决定权。

% Source PDF physical page 225; see audit/source_coverage.json.
例1，妇女同意与甲男性交，但要求使用安全套，而甲男不答应，强行发生性交，构成强奸。

例2，妇女同意与甲男性交，但要求去宾馆开房，而甲男不答应，坚持在酒吧厕所强行发生性交，构成强奸。

例3，本节开头“强行玩3P案”中，甲乙构成强奸罪。

\subsubsection*{（二）构成要件}

本罪的行为结构：

\InsertDiagram{figures/manual_p225_1.tex}

1.行为主体。单独直接实行犯只能是男子。妇女可以成为强奸罪的教唆犯、帮助犯、间接正犯和共同正犯。由于妇女可以构成强奸罪的间接正犯，因此强奸罪不是真正身份犯。

例1，甲（女）教唆乙（男）强奸丙（女），乙便实施了强奸。甲是强奸罪的教唆犯。

例2，甲（女）唆使乙（男，精神病患者）强奸丙（女），乙便实施了强奸。甲是强奸罪的间接正犯。

例3，甲（女）、乙（女）和丙（男）共谋强奸丁（女），甲乙使用暴力打晕丁，丙上前奸淫丁。甲乙丙是强奸罪的共同正犯（共同实行犯）。暴力行为属于强奸罪的实行行为。

2.实行行为。两个行为：一是暴力、胁迫或其他手段；二是奸淫行为（性交行为）。

（1）暴力手段，是指对妇女行使有形力一使妇女无法反抗。

（2）胁迫手段。行为模型：甲对乙以恶害相通告：“你若不答应我的要求，我就对你施加恶害。”\ensuremath{\rightarrow}由此使乙对甲产生恐惧心理一乙基于恐惧心理而答应了甲的不法要求。

一是“以恶害相通告”。例如，副导演甲告诉女演员乙：“跟我上床，给你女一号。”乙答应。事后甲未兑现。甲不构成强奸，因为甲是\CJKunderwave{以好处相利诱}，而非以恶害相通告。这是\CJKunderwave{强奸与权色交易的区别}。

二是“你不答应，\CJKunderwave{我}会对你施加恶害”。例如，乙女（18岁）刚来深圳，借住在甲男家里，甲要求与乙上床，否则请其搬出去。乙便准备搬出去，甲欺骗乙：“这么晚你如果出门，不出10分钟，一定会被列徒强奸！”乙很害怕，便不敢出去，答应了甲的要求。甲声称的是“外面的歹徒会对你施加恶害”，不属于胁迫乙，不构成强奸罪。

三是“你不答应，我会对\CJKunderwave{你}施加恶害”。例如，甲男与同事乙女关系要好，甲男对乙女称：“我深深爱上你了，你若不答应跟我上床，我就自杀！”乙担心甲，只好与甲上床。甲的行为属于“你不答应，我就对自已施加恶害”，不属于胁迫乙，不构成强奸罪。

四是妇女产生“恐惧心理”，基于恐惧心理，被迫与行为人发生性行为。例如，甲晚上拦路强奸，将一位女孩打倒在地，女孩忽然发现甲面容俊秀，心里嗔怪：既然能靠颜值，为何要动粗？便装作不敢反抗，与甲发生了性行为。甲构成强奸罪，但属于未遂。\CJKunderwave{注}\CJKunderwave{意}：被害人有承诺，并不要求行为人知晓。

% Source PDF physical page 226; see audit/source_coverage.json.
五是胁迫的程度与效果，既包括抢劫罪中的胁迫效果，例如，以暴力相胁迫，使妇女无法反抗，也包括敲诈勒索罪中的胁迫效果，例如，以揭发隐私相胁迫，使妇女恐惧而不敢反抗。\CJKunderwave{结论}：强奸罪中的“胁迫”=抢劫罪中的“胁迫”+敲诈勒索罪中的“胁迫”。

（3）其他手段，是指除暴力、胁迫外的其他使妇女不知反抗、不能反抗的强制手段。

第一，昏醉强奸。例如，用酒灌醉或者药物麻醉后进行强奸。\CJKunderwave{注意}：妇女自行昏醉，甲路过趁机发生性行为，构成强奸罪，这是因为妇女昏醉便失去了承诺能力，此时发生性行为视为违背妇女意愿。此时强奸罪的实行行为便只有奸淫行为，不需要强制手段。需要强制手段的前提是妇女意识清醒，此时需要通过强制手段压制妇女反抗。

第二，欺骗型强奸。例如，冒充妇女的丈夫或情夫进行强奸。\CJKunderwave{注意}：事实错误会导致妇女的承诺无效，但动机错误不影响妇女的承诺。例如，欺骗女演员“陪睡，就给女一号”，不构成强奸罪。对此可参见总论“被害人承诺”一节。

3.行为对象是妇女。针对已满14周岁的妇女，是普通强奸；针对不满14周岁的幼女，是奸淫幼女。

（1）婚内强奸问题。从第236条的规定看，并没有将婚内强奸排除在强奸罪之外。丈夫强奸妻子符合强奸罪的构成要件。所谓婚内强奸是否构成强奸罪的问题，主要是考虑证据证明、刑事政策等问题。（2015年第51题）

（2）奸淫幼女。

\InsertDiagram{figures/manual_table_p226_1.tex}

（3）精神病女和痴呆女。第一，这类妇女对性行为缺乏理解能力，缺乏承诺能力。因此，明知对方是这类妇女，即使其“同意”，与其发生性行为的，也构成强奸罪。（2014年第51题）。第二，为了与这类妇女形成事实婚姻关系而与其共同生活，在共同生活期间发生性行为的，不应认定为强奸罪。

\subsubsection*{（三）法定刑升格条件}

1．“在公共场所当众强奸妇女、奸淫幼女”

（1）“公共场所”，是指多数人或不特定人可以进出的场所。

（2）“当众强奸”是指明知能够为多数人或不特定人知晓，仍实施强奸。这里的“众”不要求必须有3人在场，主要指有不特定人在场。这里的“众”不包括共犯人。

2.“二人以上轮奸”

（1）轮奸，是指对连续两次强奸的违法事实负正犯责任。

例1（双方都构成轮奸），甲乙共谋强奸丙女，甲先强奸，乙后强奸。甲乙构成轮奸。

例2（单方构成轮奸），甲打晕丙女，强奸了丙女，又指使路过的乙强奸了丙女。乙

% Source PDF physical page 227; see audit/source_coverage.json.
不构成轮奸，甲构成轮奸，甲要对乙的强奸罪负共同正犯的责任。\footnote{甲趁丙女熟睡，强奸了丙女。丙女醒来后，甲教唆乙强奸丙女，乙照办。甲对乙的强奸罪只负教唆犯的责任，因此甲不构成轮奸。}又如，甲打晕丙女，由乙强奸。乙离开后，甲又强奸了丙女。乙不知道甲强奸了丙女。乙不构成轮奸，甲构成轮奸。

（2）轮奸未遂。轮奸是情节加重犯，在成立的基础上，也有既遂、未遂、中止等问题。若两个强奸犯自身均强奸既遂，才能构成轮奸既遂。

例如，甲乙共谋轮奸丙女，共同将丙打晕，甲先奸淫得逞，乙因自身原因未能得逞或心生怜悯而放弃。就基本犯而言，甲构成强奸罪既遂，基于“部分实行、全部负责”原则，乙应对甲的既遂结果负责，就此而言乙也构成强奸罪既遂。就情节加重犯而言，甲乙成立轮奸。由于在事实层面，乙自身没有强奸既遂，所以轮奸属于未遂。所以，甲构成轮奸未遂，乙构成轮奸未遂或中止。甲同时触犯“强奸罪既遂”和“轮奸未遂”，想象竞合，择一重刑论处。乙同时触犯“强奸罪既遂”和“轮奸未遂或中止”，想象竞合，择一重刑论处。（2024年试题）

3.“奸淫不满10周岁的幼女或者造成幼女伤害”

奸淫不满14周岁的幼女，从重处罚。奸淫不满10周岁的幼女，加重处罚（适用升格法定刑)。

4.\CJKunderwave{“致使被害人重伤、死亡或者造成其他严重后果”，其中蕴含结果加重犯（强奸罪}\CJKunderwave{致人重伤、死亡）}

（1）“因”是实行行为。必须是强奸罪的实行行为导致重伤或死亡结果。所谓实行行为，是指带着奸淫目的实施的暴力和奸淫行为。不能评价为强奸罪的实行行为的其他行为导致重伤或死亡结果，不属于强奸致人重伤或死亡。对其他行为导致重伤、死亡需独立评价。

例1，甲强奸妇女后，为了灭口杀死妇女，定强奸罪和故意杀人罪，并罚。

例2，甲强奸妇女，妇女反抗中咬断甲的舌头，甲恼羞成怒，为了泄愤杀死妇女。这里杀死妇女的行为已经不能评价为强奸行为，而是独立的故意杀人。对甲定强奸罪（未遂）和故意杀人罪（既遂），并罚。（2007年第12题）

（2）因果关系。强奸行为与加重结果之间应有直接因果关系。

例如，甲强奸妇女后离开现场，妇女昏迷，苏醒后爬两步，不慎掉入湖中淹死。甲构成强奸致人死亡。

（3）主观心理。由于强奸行为包括暴力行为，使用暴力致人重伤、死亡，完全可能是故意所为。因此，强奸致人重伤、死亡，包括故意致人重伤、死亡。例如，甲欲强奸乙女，不顾死活地将乙打成重伤昏迷，然后奸淫，后乙因伤势过重死亡。

（4）“强奸致人重伤、死亡”中的“人”只包括被害妇女，不包括前来阻挡的第三人。例如，甲强奸妇女，乙前来阻挡，甲为了排除障碍，将乙打成重伤，进而顺利强奸了妇女。甲构成强奸罪和故意伤害罪，并罚。

（5）“造成其他严重后果”包括：导致被害人自杀、精神失常、怀孕（依据典型判

例）。注意：被害人自杀，行为人不属于“强奸罪致人死亡”这一结果加重犯，因为缺乏刑法上的因果关系。

% Source PDF physical page 228; see audit/source_coverage.json.
\subsubsection*{（四）罪数问题}

注意：刑法只在拐卖妇女、儿童罪中规定了结合犯：拐卖罪+强奸罪=拐卖罪（加重处罚），在其他罪中又犯强奸罪，一律按原则办事，数罪并罚。（2007年第12题）

\subsection*{二、负有照护职责人员性侵罪}

第236条之一对已满十四周岁不满十六周岁的未成年女性负有监护、收养、看护、教育、医疗等特殊职责的人员，与该未成年女性发生性关系的，处三年以下有期徒刑；情节恶劣的，处三年以上十年以下有期徒刑。

有前款行为，同时又构成本法第二百三十六条规定之罪的，依照处罚较重的规定定罪处罚。[《刑法修正案（十一）》增设]

\subsubsection*{（一）构成要件}

1.行为主体：负有照护职责的人。

2.行为对象：已满14周岁不满16周岁的少女。

3.行为方式：发生性交行为，不包括猥亵行为。

4.被害人承诺无效。

\InsertDiagram{figures/manual_p228_1.tex}

\subsubsection*{（二）司法认定问题}

1.行为主体的身份

不能认为，只要是老师、医生，就一定是本罪的行为主体。只有负有较稳定的照护职责（照管、看护职责）的人才能成为本罪的行为主体。

例1，15周岁的少女乙在课外补习班上课期间，自愿与补习老师甲男发生性关系。甲不构成本罪。

例2，狱警甲负责监管15周岁的乙女，乙女同意与甲发生性行为。甲构成本罪。

例3，乙需要出差三天，因对15周岁的女儿独自在家不放心，请邻居甲男有空照顾一下女儿。期间，该少女与甲自愿发生性关系。甲不构成本罪。

2.行为对象的年龄

不满16周岁是真正的年龄条件。已满14周岁不是真正的年龄条件，只是用于与“奸淫幼女型”强奸罪区分的界限要素。

负有照护职责的人在未满16周岁女孩同意的情况下，与其发生性交行为，构成犯罪，

% Source PDF physical page 229; see audit/source_coverage.json.
只是罪名有所不同。当女孩已满14周岁时，罪名是负有照护职责人员性侵罪，是轻罪；当女孩未满14周岁时，罪名是强奸罪，是重罪。两个罪名是性质相同、程度不同的关系，重罪能够包容评价为轻罪。当在两个罪名之间产生认识错误时，在轻罪的范围内主客观相一致，定轻罪。

例如，乙女13周岁，长期住院治疗，负责医生是甲男。乙对甲谎称自己15周岁，甲信以为真。经乙同意，甲与乙发生性交行为。主观上，甲以为自己犯负有照护职责人员性侵罪；客观上，甲属于奸淫幼女型强奸。后罪能够包容评价为前罪。甲构成前罪。

\subsection*{三、强制猥亵、侮辱罪}

第237条第1、2款以暴力、胁迫或者其他方法强制猥亵他人或者侮辱妇女的，处五年以下有期徒刑或者拘役。

聚众或者在公共场所当众犯前款罪的，或者有其他恶劣情节的，处五年以上有期徒刑。[本条经《刑法修正案（九）》修正］

本罪侵犯的法益是他人的性自主权。这里的性自主权主要表现为性羞耻心。

1.行为主体是已满16周岁的自然人。

2.行为对象是已满14周岁的人，包括男性和女性。

（1）如果猥亵不满14周岁的儿童，构成猥亵儿童罪。

（2）如果强奸女性，包括不满14周岁的幼女和已满14周岁的妇女，均构成强奸罪。如果强奸不满14周岁的男童，构成猥亵儿童罪；如果强奸已满14周岁的男性，由于我国强奸罪的对象尚不包括男性，只能以强制猥亵罪论处。

3.猥亵、侮辱

（1）猥亵，是指侵犯性羞耻心的行为，主要情形：一是直接对他人实施猥亵行为；二是强迫他人对行为人自己或者第三人实施猥亵；三是强迫他人自行猥亵；四是强迫他人观看其他人猥亵。\CJKunderwave{注意}：不包括偷看他人洗澡、做爱，不包括在公共场所露阴、手淫但不强迫他人观看。

（2）侮辱。这里的“侮辱”的含义与猥亵相同，也是侵犯性羞耻心的行为。\CJKunderwave{注意}：这里的“侮辱”与侮辱罪中的“侮辱”有所区别，后者侵犯的仅仅是他人的名誉。（2013年第16题，2016年第58题）

4.强制方法。必须以暴力、胁迫或者其他强制方法猥亵、侮辱。强制方法的特点是使他人不能反抗、不敢反抗、不知反抗。例如，甲强迫乙与自已裸聊、裸露直播，构成本罪。

[注意］强迫发裸照问题。例如，甲强迫乙给自已发裸照、裸露视频，乙拍裸照、视频时是关起门来的。结论：（1）乙若是成年人，则甲不构成本罪，因为强制猥亵的对象应是裸露的身体，而不包括物品。（2）乙若是年满14未满18周岁，则甲构成本罪。这是司法解释对未成年人的特别保护。\footnote{2023年6月1日最高人民法院、最高人民检察院《关于办理强奸、猥亵未成年人刑事案件适用法律若干问题的解释》。}

5.强制猥亵、侮辱罪与强奸罪不是A与-A的对立排斥关系，而是A与A+B的包容评价关系。例如，甲男强扒妇女衣服时被抓。无法证明甲的主观故意是强奸还是猥亵。此时，可认定甲构成强制猥亵罪。

% Source PDF physical page 230; see audit/source_coverage.json.
\subsection*{四、猥亵儿童罪}

第237条第3款猥亵儿童的，处五年以下有期徒刑；有下列情形之一的，处五年以上有期徒刑：

\subsubsection*{（一）猥亵儿童多人或者多次的；}

\subsubsection*{（二）聚众猥亵儿童的，或者在公共场所当众猥亵儿童，情节恶劣的；}

\subsubsection*{（三）造成儿童伤害或者其他严重后果的；}

\subsubsection*{（四）猥亵手段恶劣或者有其他恶劣情节的。［《刑法修正案（十一）》修订]}

1.行为主体：已满16周岁的人，包括男性和女性。主要情形包括：男性对男童、女童猥亵，女性对男童、女童猥亵。

2.行为对象。（1）猥亵女童，不包括性交行为，否则构成强奸罪。（2）猥亵男童，包括性交。

3.被害人承诺无效。即使儿童自愿，行为人也构成本罪。

4.行为方式。让儿童给自己发裸照、裸露视频，构成本罪。理由是，本罪的保护法益不是性羞耻心，而是儿童的身心健康。该行为方式侵害了儿童的身心健康。\footnote{2023年6月1日最高人民法院、最高人民检察院《关于办理强奸、猥亵未成年人刑事案件适用法律若干问题的解释》。}

[总结］强奸罪、强制猥亵罪、猥亵儿童罪

\InsertDiagram{figures/manual_table_p230_1.tex}

\section{侵犯人身自由的犯罪}

[问题】机长在能降落飞机的情况下，欺骗乘客“天气不好，无法降落”，故意不降落。机长是否属于非法拘禁？（故意不降落案）

\subsection*{一、非法拘禁罪}

第238条（第1款第1句）非法拘禁他人或者以其他方法非法剥夺他人人身自由的，处三年以下有期徒刑、拘役、管制或者剥夺政治权利。（第1款第2句）具有殴打、侮辱情节的，从重处罚。

（第2款第1句）\CJKunderwave{犯前款罪，致人重伤的}，处三年以上十年以下有期徒刑；致人死亡的，处十年以上有期徒刑。（第2款第2句）\CJKunderwave{使用暴力致人伤残、死亡的}，依照本法第二百三十四条、第二百三十二条的规定\CJKunderwave{（故意伤害罪、故意杀人罪）定罪处罚}。

（第3款）\CJKunderwave{为索取债务非法扣押、拘禁他人的}，依照前两款的规定处罚。

% Source PDF physical page 231; see audit/source_coverage.json.
（第4款）\CJKunderwave{国家机关工作人员}利用职权犯前三款罪的，依照前三款的规定从重处罚。

\subsubsection*{（一）构成要件}

拘禁行为：非法剥夺他人人身自由的行为。\footnote{该罪是继续犯。按照司法解释规定，非法拘禁持续24小时，才构成既遂。未满24小时，就是未遂。这种未遂一般不值得定罪处罚。不过，如果又有其他严重情节，则对这种未遂值得处罚。}

例1，向警察报假案，欺骗警察刑事拘留了他人，构成非法拘禁罪的间接正犯（同时也触犯诬告陷害罪）。

例2，驾驶汽车高速行驶，使欲下车的乘客基于恐惧心理无法下车。

例3，发现将同学错关在房间里，却故意不开门，构成不作为的非法拘禁罪。

[注意］欺骗问题

例1，甲欺骗乙进入电梯后关闭电源，导致乙出不去。甲对乙谎称停电检修，请乙耐心等待。乙相信并照办。甲的欺骗行为导致乙认为自己只能待在电梯里，失去了人身自由。甲构成非法拘禁罪。

[结论一］欺骗行为导致被害人认为自己出不去了，行为人构成非法拘禁罪。

例2，本节开头“故意不降落案”中，机长欺骗乘客，乘客认为自已出不去了。机长构成非法拘禁罪。

例3，甲欺骗乙：“为了给你做体检，你需要在这个房间待上一段时间，房门没锁，但请不要出去”。乙相信并照办。乙认为自已能出去，只是为了体检不想出去。甲不构成非法拘禁罪。\footnote{如果甲背着乙向乙的亲人勒索财物，也不构成绑架罪，因为成立绑架罪，要求对人质非法拘禁。甲构成诈骗罪和敲诈勒索罪，想象竞合。}

[结论二］欺骗并未导致被害人认为自己出不去了，只是基于某种动机不想出去，行为人不构成非法拘禁罪。

例4，甲欲强奸乙女，谎称开车送乙回家，实际是想开到郊外。乙不知情而上车，行驶一段路程后，乙发现路线不对，要求停车，甲便停车，乙下车。甲构成强奸罪预备阶段的中止。乙坐在车上时，并未认为自已下不了车，只是想回家而没下车。因此，甲不构成非法拘禁罪。如果甲不让乙下车，则构成非法拘禁罪。

\subsubsection*{（二）第2款第1句：“致人重伤、死亡”}

这是本罪的结果加重犯，也即非法拘禁罪（过失）致人重伤、死亡。根据结果加重犯的基本原理，注意三点：

1．“因”的要求。导致加重结果的“因”是本罪的实行行为。判断本罪的实行行为，可根据“行为与故意或目的同时存在原则”。

例如，甲非法拘禁乙两天，已经构成非法拘禁罪；第三天，甲忽然发现乙死亡，原来甲不慎将绳子勒得太紧，导致乙室息死亡。死亡结果是绳子勒的行为，这是非法拘禁罪的实行行为，因此构成非法拘禁罪致人死亡。

2.“因”与“果”要有因果关系。在此注意介入因素。（2014年第59题，2016年第58题）

例1，甲为了索债，将乙关在屋子里，甲外出。乙感觉无力还债，便跳楼自杀身亡。

% Source PDF physical page 232; see audit/source_coverage.json.
该介入因素异常，死亡结果与甲的拘禁行为没有刑法上的因果关系。甲不构成非法拘禁罪致人死亡。（2015年第8题）

例2，甲非法拘禁乙，乙悄悄解开绳子逃跑，甲发现后紧追，乙跑向阳台呼救时不慎跌落摔死。该介入因素不异常，乙的死亡与甲的拘禁行为有因果关系，甲构成非法拘禁罪致人死亡。

3.主观上，对加重结果持过失心理。如果持故意心理，则不属于非法拘禁罪致人重伤、死亡。例如（2012年第16题），甲为了非法拘禁乙，大力反扭乙的胳膊，致其胳膊折断（重伤）。此时伤害行为是为了拘禁乙，也是非法拘禁罪的实行行为。甲的一个伤害行为同时触犯非法拘禁罪和故意伤害罪（重伤），想象竞合，择一重罪论处，定故意伤害罪（重伤）。

\subsubsection*{（三）第2款第2句：“使用暴力致人伤残、死亡的，应以故意伤害罪、故意杀人罪论处}

注意，第2款第1句的结果加重犯是一个行为。而第2款第2句的暴力行为是在非法拘禁罪的行为之外的第二个行为。

第一种观点（多数说）：条文中的“致人伤残、死亡”是指过失致人伤残、死亡，但拟制为故意伤害罪、故意杀人罪。这是一个法律拟制性规定（特事特办），也被称为转化犯。

第二种观点（少数说）：条文中的“致人伤残、死亡”是指故意致人伤残、死亡。所以，“以故意伤害罪、故意杀人罪论处”，这是注意规定。（正常办事）（2025年主观题）

按照第一种观点继续讲解：

1.“因”的要求。导致伤残、死亡结果的“因”不能是非法拘禁罪的实行行为，否则就构成结果加重犯“非法拘禁罪致人重伤、死亡”。这里的“因”只能是非法拘禁罪的实行行为之外的暴力行为，属于第二个行为。这是二者的关键区分。对此判断，可根据“行为与故意或目的同时存在原则”。（2011年第60题）

例如，甲非法拘禁了乙，次日，乙辱骂甲，甲想教训乙，扇耳光，碰巧打到乙的太阳穴，导致乙死亡。甲扇耳光，目的不是拘禁乙，而是泄愤，因此不是非法拘禁罪的实行行为，该暴力行为过失致人死亡，应拟制为故意杀人罪。

在罪数上，对甲最终只定故意杀人罪，不再处罚非法拘禁罪。理由是，刑法将过失致人死亡拟制为故意杀人罪，是有背景条件的，例如，在非法拘禁罪的背景下。刑法不可能将单纯的过失致人死亡拟制为故意杀人罪，否则过失致人死亡罪便没有存在必要了。所以，拟制的公式是：非法拘禁罪+过失致人死亡=故意杀人罪。由于拟制为故意杀人罪已经用到了（评价到了）非法拘禁罪，所以定了故意杀人罪后，不能再与非法拘禁罪数罪并罚，否则就意味着将非法拘禁罪评价了两次（处罚了两次），违反了禁止重复评价（处罚）原则。

2.主观要求。这里的“致人伤残、死亡”，是指过失所为，被拟制为故意伤害罪、故意杀人罪。如果故意将人打成重伤或故意杀人，则是正常的故意伤害罪（重伤）、故意杀人罪。

例如，甲对乙已经构成非法拘禁罪后，乙辱骂甲，甲生气，便杀了乙。甲构成非法拘禁罪和故意杀人罪。在罪数上，少数说认为，故意杀人罪吸收非法拘禁罪。理由是，两罪都是人身犯罪，属于性质相同、程度不同的关系，重罪可以吸收轻罪。多数说认为，两罪应并罚。理由是，虽然两罪都是人身犯罪，但是人身自由和生命是性质不同的法益。如果只定故

% Source PDF physical page 233; see audit/source_coverage.json.
意杀人罪，则会遗漏评价非法拘禁罪制造的法益侵害事实，违反了禁止遗漏评价原则。（2021年主观题、2022年主观题）

\InsertDiagram{figures/manual_table_p233_1.tex}

\subsection*{二、绑架罪}

第239条以勒索财物为目的绑架他人的，或者绑架他人作为人质的，处十年以上有期徒刑或者无期徒刑，并处罚金或者没收财产；情节较轻的，处五年以上十年以下有期徒刑，并处罚金。

犯前款罪，杀害被绑架人的，或者故意伤害被绑架人，致人重伤、死亡的，处无期徒刑或者死刑，并处没收财产。

以勒索财物为目的偷盗婴幼儿的，依照前两款的规定处罚。[本条第2款经《刑法修正案（九）》修正]

绑架罪，是指带着胁迫第三人的目的，用强制手段实力控制人质的行为。

% Source PDF physical page 234; see audit/source_coverage.json.
\InsertDiagram{figures/manual_p234_1.tex}

\subsubsection*{（一）构成要件}

1.行为一与行为二

例如，甲欲向B勒索财物，绑架了B的老婆A，实力控制了A，尚未开始向B勒索财物时，甲被抓。

第一种观点（少数说）：甲构成绑架罪未遂。理由是，绑架罪的保护法益既包括人质的人身法益，也包括被勒索者的财产法益。只有实施了行为二（向B勒索财物的行为），才构成绑架罪既遂。绑架罪是复合行为犯，实行行为有两个：绑架行为+勒索行为。

第二种观点（多数说）：甲构成绑架罪既遂。理由是，绑架罪的保护法益只有人质的人身法益。绑架罪是单一行为犯，实行行为只有绑架行为。

按照多数说，结论：成立绑架罪，只要求实施行为一，不要求实施行为二（向B勒索财物）。但在实施行为一时，必须具有目的二（向B勒索财物的目的），否则就只是非法拘禁罪。

2.目的二

绑架罪的目的二，是指向第三人提出不法要求的目的，也即胁迫、威胁第三人的目的。常见的不法要求是勒索财物，也有其他不法要求。例如，甲绑架了一个过路妇女，要求警方释放囚犯乙（甲的哥哥），否则就杀死该妇女。甲成立绑架罪。

[总结]\CJKunderwave{绑架罪与相似罪的区分}。

\CJKunderwave{绑架罪与相似罪名的客观行为可以相同}，\CJKunderwave{都可以是实力控制行为}。\CJKunderwave{关键区分在于有无}\CJKunderwave{目的二（向第三人勒索财物的目的）}。

（1）绑架罪与非法拘禁罪的区分：前者有目的二，后者无目的二。非法拘禁罪是A，绑架罪是A+B，B是目的二。绑架罪可以包容评价为非法拘禁罪。触犯绑架罪必然触犯非法拘禁罪，二者是法条竞合关系。（2013年第56题）

（2）绑架罪与拐卖妇女、儿童罪的区分：前者的目的二是向第三人勒索财物，后者的目的是出卖。例如，狗蛋带着向小芳丈夫要钱的目的，拘禁了小芳，构成绑架罪。又如，狗蛋带着出卖小芳的目的，拘禁了小芳，构成拐卖妇女罪。

（3）绑架罪与抢劫罪的区分：前者的目的二是向第三人勒索财物，后者的目的是向被绑的人勒索财物。例如，狗蛋带着向小芳要钱的目的，拘禁了小芳，构成抢劫罪。

3.总结成立条件

(1）带着目的二，也即带着向第三人提出不法要求的目的（主要是勒索财物的目的)。

% Source PDF physical page 235; see audit/source_coverage.json.
（2）对合格的人质实施实力控制行为。

（3）合格的人质是：从行为时、一般人的角度看，会使第三人产生担忧的人。

例1，狗蛋绑架了柏浪涛，问普京要钱。普京根本不会担忧柏浪涛的安危，因此狗蛋不构成绑架罪，而仅构成非法拘禁罪。

例2（2021年试题），甲绑架了乙，要求乙的妻子丙交付赎金。丙心想，乙经常打骂自己，便不给钱。从行为时、一般人的角度看，绑架丈夫，会使妻子产生担忧，所以甲成立绑架罪，而且既遂。也即，合格的人质，只要求有使第三人担忧的一般可能性，不要求使第三人实际担忧。

4.既遂条件：实际控制了合格的人质。不要求实施行为二，不要求要到钱。（2011年第54题）

（1）\CJKunderwave{实际控制了}人质。如果没有实际控制，则是未遂。例如，甲带着向乙要钱的目的，着手绑架乙的孩子丙，丙奋力反抗，逃脱。甲构成绑架罪未遂。

（2）实际控制了\CJKunderwave{人质}。如果实际控制的不是合格的人质，则是未遂。

例1（2017年主观题），甲带着向乙要钱的目的，着手绑架乙的孩子丙，守在丙的小学门口，丙与同学丁出来，甲动手劫持丙时，误以为丙不是乙的孩子，误以为丁是乙的孩子，将丁劫持走，然后向乙打电话要钱。乙看到丙在家里，骂了句“骗子”，便挂了电话。甲没有实力控制到合格的人质，构成未遂。但甲构成非法拘禁罪既遂，因为非法拘禁罪的对象是一般对象，不要求是人质。甲同时构成绑架罪未遂和非法拘禁罪既遂，想象竟合，择一重罪论处。\footnote{可能有人认为，甲构成对象错误，应构成绑架罪既遂。对此，请参见总论“主观要件”一讲中关于对象错误的分析。}

例2，甲带着向小芳要钱的目的，着手绑架小芳的孩子，守在一个小学门口，看到小狗蛋出来，误以为是小芳的孩子，予以劫持，然后向小芳打电话要钱。小芳骂道：“人家一直单身，就没孩子，骗子！”原来甲情报有误。甲不构成绑架罪，仅构成非法拘禁罪。注意：例2与例1的区别在于，例1中，甲有绑架到合格人质的可能性、危险性，有让乙担忧的可能性。而例2中没有这种可能性、危险性。

[总结］第一，成立绑架罪，要求有绑架到合格人质的可能性、危险性。第二，成立绑架罪既遂，要求绑架到合格人质。\footnote{绑匪将合格人质绑到手就既遂，既遂后，并向第三人勒索到了钱财，又构成敲诈勒索罪既遂。对绑架罪既遂和敲诈勒索罪既遂不并罚，二者是手段与目的的牵连犯关系，择一重罪论处，定绑架罪既遂。}

\subsubsection*{（二）第238条第3款的解读}

该款规定：“为索取债务非法扣押、拘禁他人的，以非法拘禁罪论处。”例如，甲向乙讨债，乙不还，甲便将乙的妻子非法扣押，逼迫乙还债。甲构成非法拘禁罪。（2025年主观题）

1.条文属性：注意规定而非法律拟制。也即，这种情形原本就要定非法拘禁罪，而非原本定绑架罪，例外地按非法拘禁罪论处。这是因为，成立绑架罪要求行为人具有目的二，也即向第三人勒索财物的目的，这是一种非法占有他人财物的目的。而条文的情形是行为人为了索取债务，为了主张自己的债权，不具有非法占有目的，所以不构成绑架罪，只构成非法拘禁罪。

% Source PDF physical page 236; see audit/source_coverage.json.
2.行为方式有三种情形：

（1）拘禁债务人亲属，向债务人索债。

（2）拘禁债务人，向其亲属索债。这里的亲属仅包括具有共同财产关系的亲属。如果向除此之外的亲友索债，就表明具有非法占有该亲友财物的目的，就构成绑架罪。例1，甲向乙讨债，乙不还。甲将乙非法扣押，逼迫乙的妻子还债。甲构成非法拘禁罪。例2，甲向乙讨债，乙不还。甲将乙非法扣押，逼迫乙的表弟替表哥还债。甲构成绑架罪。

（3）拘禁债务人，向债务人本人索债。例如，甲向乙讨债，乙不还。甲将乙非法扣押，逼迫乙还债。这类似于抢劫，但由于没有非法占有目的，故不构成抢劫罪，只构成非法拘禁罪。（2014年第59题）

3.“债务”的认定

（1）司法解释对“债务”作了扩大解释，包括赌债、高利贷等非法债务。

（2）如果形式上是索取债务，但没有任何事实根据，实质上具有非法占有目的，则构成绑架罪。

例1，甲伪造借条，声称乙欠自已钱财，非法拘禁了乙的妻子，威胁乙还债。索取这种没有事实根据的“债务”，显然具有非法占有目的，构成绑架罪。

例2，乙欠甲10万元，甲非法拘禁了乙的妻子，威胁乙还100万元。甲具有非法占有目的，构成绑架罪。

\subsubsection*{（三）第239 条第2款的解读}

该款规定：“犯前款罪，杀害被绑架人的，或者故意伤害被绑架人，致人重伤、死亡的，处无期徒刑或者死刑，并处没收财产”

1.绑架+杀害

这是结合犯：绑架罪+故意杀人罪=绑架罪（加重处罚）。（2014年第16题）

（1）杀害结果：要求杀死人质（多数说）。例如，甲绑架了乙后，向乙的老公要钱，没要到，在杀乙时，被警察抓捕。不能适用该结合犯，而应正常处理，也即，绑架罪（既遂）与故意杀人罪（未遂），数罪并罚。

（2）杀害时间：从绑架的着手到释放人质前。

例1，甲使用暴力要绑架乙，但无法控制乙，干脆杀死乙。这属于，绑架罪（未遂）+杀人（既遂）=绑架罪（加重处罚）。

例2，甲勒索到赎金后，释放了人质乙，又后悔，开车追了三公里，追上人质，将其撞死。此时不属于“杀害被绑架人”，因为乙已经恢复自由，不属于被绑架人（人质）。此时杀人行为要与绑架罪数罪并罚。（2020年试题）

2.绑架+伤害

行为公式：绑架罪+故意伤害罪致人重伤、死亡=绑架罪（加重处罚）。

（1）“故意伤害罪致人重伤、死亡。”第一，不包括轻伤。第二，对重伤是故意，对死亡是过失。所以，要构成这里的故意伤害罪，要求：客观上具有造成重伤的可能性，主观上具有造成重伤的故意（包括间接故意）。

（2）伤害时间：从绑架的着手到释放人质前。

% Source PDF physical page 237; see audit/source_coverage.json.
（3）伤害行为的性质。注意，直接故意的杀害行为不能成为绑架罪的实行行为，因为绑架罪只是控制人质的人身自由；如果直接想杀死人，就不构成绑架罪，因为活人才能成为人质。\footnote{间接故意的杀害行为有可能成为绑架罪或非法拘禁罪的实行行为。例如，甲为了向乙的家人丙勒索财物，欲绑架乙，乙激烈反抗，甲为了制服乙，使用暴力，明知该暴力行为可能导致死亡，但甲带着放任心理实施暴力，导致乙重伤，制服了乙，甲向丙要钱。等丙答复期间，乙因伤势过重而死亡。甲属于一个行为同时构成绑架罪和故意杀人罪既遂，本应按照想象竞合择一重罪论处，但是法条有特殊规定“绑架罪+故意杀人罪=绑架罪（加重处罚）”，那么就按照该法条规定处理。不过，该处理不属于结合犯，因为结合犯是两个行为，而属于结果加重犯，也即“绑架罪（间接故意）致人死亡”。}但是，为了控制目的而实施的伤害能够成为绑架罪的实行行为。所以，在绑架罪中，伤害行为有两种情形。

第一，伤害行为不是绑架罪的实行行为。例如，甲绑架了乙，没能向乙的家人勒索到钱，很生气，将乙打成重伤。这属于“绑架罪+故意伤害罪致人重伤、死亡=绑架罪（加重处罚）”。此时是结合犯的性质。

第二，伤害行为是绑架罪的实行行为。例如，甲为绑架乙，欲实力控制乙，乙反抗，甲故意将乙打成重伤，制服了乙，向乙的家人勒索赎金。这也属于“绑架罪+故意伤害罪致人重伤、死亡=绑架罪（加重处罚）”。此时是结果加重犯的性质，可表述为“绑架罪（故意）致人重伤”。

［巩固练习］题1，甲用绳子绑住乙，向乙的家人要钱，未要到钱，回来发现乙死了，原来绳子绑得太紧，导致乙室息死亡。“用绳子绑”的目的是控制住乙，目的不是将乙弄成重伤，因此该行为不构成故意伤害罪，便不构成故意伤害致人死亡，而构成过失致人死亡罪，因此不属于“绑架罪+故意伤害罪致人重伤、死亡=绑架罪（加重处罚）”。甲的一个实力控制行为同时触犯绑架罪与过失致人死亡罪，想象竞合，择一重罪论处，定绑架罪。\footnote{旧条文规定了结果加重犯“绑架罪（过失）致人死亡”。这种情形便属于“绑架罪（过失）致人死亡”。《刑法修正案（九）》删掉了结果加重犯的规定。由此，结果加重犯还原为想象竞合犯。}

题2，甲用绳子绑住乙，乙呼救，甲用毛巾塞住乙的嘴巴。甲向乙的家人要钱，未要到钱，回来发现乙死了，原来是毛巾导致乙室息死亡。“用毛巾塞住嘴巴”的目的是控制住乙，仍属于控制乙的一种辅助手段。其目的不是将乙弄成重伤，因此该行为不构成故意伤害罪，便不构成故意伤害致人死亡，而构成过失致人死亡罪，因此不属于“绑架罪+故意伤害罪致人重伤、死亡=绑架罪（加重处罚）”。甲的一个实力控制行为同时触犯绑架罪与过失致人死亡罪，想象竞合，择一重罪论处，定绑架罪。（2020年试题）

\InsertDiagram{figures/manual_table_p237_1.tex}

% Source PDF physical page 238; see audit/source_coverage.json.
\ContinuedTable{figures/manual_table_p238_1.tex}

\subsubsection*{（四）共犯模型}

模型1，甲绑架了丙，向丙的妻子勒索钱财，期间，让乙帮忙为丙做饭。乙照办。乙不构成绑架罪的承继的共同犯罪，因为做饭不是犯罪行为，没有侵害法益。

模型2，甲绑架了丙，向丙的妻子勒索钱财，期间，让乙看押丙。乙照办。乙构成绑架罪的承继的共同犯罪，因为绑架罪是继续犯。

模型3，甲绑架了丙，欲向其妻子勒索财物，欺骗乙：“丙的妻子欠我钱，你帮我看押丙，我要到钱，会感谢你。”乙信以为真而照办。甲乙在非法拘禁罪的范围内成立共同犯罪，对甲定绑架罪，对乙定非法拘禁罪。（2016年第15题）

模型4，甲绑架了丙，向丙的妻子没能勒索到钱，要撕票，叫来乙，一起杀死丙。乙不构成绑架罪的承继的共同犯罪，因为乙没有帮助实施绑架罪的实行行为（实力控制人质）。甲乙在故意杀人罪的范围内成立共同犯罪，对甲定绑架罪（适用升格法定刑），对乙定故意杀人罪。

\OriginalPanel{典型真题}

甲为勒索财物，打算绑架富商之子吴某（5岁）。甲欺骗乙、丙说：“富商欠我100万元不还，你们帮我扣押其子，成功后给你们每人10万元。”乙、丙将吴某扣押，但甲无法联系上富商，未能进行勒索。三天后，甲让乙、丙将吴某释放。吴某一人在回家路上溺水身亡。关于本案，下列哪一选项是正确的？（2016年·卷二·15题）\footnote{［答案]甲、乙、丙构成非法拘禁罪的共同犯罪，对甲以绑架罪既遂论处，对乙、丙以非法拘禁罪既遂论处。吴某的死亡与甲、乙、丙的拘禁行为没有因果关系。本题答案：D。}

A.甲、乙、丙构成绑架罪的共同犯罪，但对乙、丙只能适用非法拘禁罪的法定刑

B.甲未能实施勒索行为，属绑架未遂；甲主动让乙、丙放人，属绑架中止

C.吴某的死亡结果应归责于甲的行为，甲成立绑架致人死亡的结果加重犯

D.不管甲是绑架未遂、绑架中止还是绑架既遂，乙、丙均成立犯罪既遂

% Source PDF physical page 239; see audit/source_coverage.json.
\subsection*{三、拐卖妇女、儿童罪}

第240条拐卖妇女、儿童的，处五年以上十年以下有期徒刑，并处罚金；有下列情形之一的，处十年以上有期徒刑或者无期徒刑，并处罚金或者没收财产；情节特别严重的，处死刑，并处没收财产：

\subsubsection*{（一）拐卖妇女、儿童集团的首要分子；}

\subsubsection*{（二）拐卖妇女、儿童三人以上的；}

\subsubsection*{（三）\CJKunderwave{奸淫}被拐卖的妇女的；}

\subsubsection*{（四）诱骗、\CJKunderwave{强迫}被拐卖的妇女卖淫或者将被拐卖的妇女卖给他人迫使其卖淫的；}

\subsubsection*{（五）以\CJKunderwave{出卖}为目的，使用暴力、胁迫或者\CJKunderwave{麻醉}方法绑架妇女、儿童的；}

\subsubsection*{（六）以\CJKunderwave{出卖}为目的，\CJKunderwave{偷盗婴幼儿}的；}

\subsubsection*{（七）造成被拐卖的妇女、儿童或者其亲属\CJKunderwave{重伤、死亡}或者其他严重后果的；}

\subsubsection*{（八）将妇女、儿童卖往境外的。}

\CJKunderwave{拐卖妇女、儿童}是指\CJKunderwave{以出卖为目的}，有\CJKunderwave{拐骗}、\CJKunderwave{绑架}、\CJKunderwave{收买}、\CJKunderwave{贩卖}、接送、中转妇女、儿童的行为之一的。

本罪的保护法益是妇女、儿童的人身自由。注意（多数说）：成年妇女如果同意放弃该法益，则行为人不构成拐卖妇女罪。例如，甲对成年妇女乙讲：“我给你介绍个有钱人家，你愿不愿意？”乙愿意。甲将乙介绍给丙家，收了丙10万元。甲不构成拐卖妇女罪。（2021年试题）

\subsubsection*{（一）实行行为}

1.实行行为一：绑架。

这里的绑架是指实力控制行为，包括偷盗婴儿。就绑架而言，本罪的行为结构：

\InsertDiagram{figures/manual_p239_1.tex}

主观上，要求行为人具有两个目的，目的一是实力控制妇女、儿童的目的，目的二是出卖目的。但是，客观上，只要求实施行为一，不要求实施行为二。只有行为一是本罪的实行行为。结论：实施了实力控制行为，就成立本罪；将妇女、儿童控制到手（拐到手），就构成既遂。（2007年第55题，2010年第61题，2021年试题）

[总结]\CJKunderwave{本罪与他罪的区分关键就在于目的二}。

（1）本罪与非法拘禁罪的区分：前者有目的二，后者无目的二。

（2）本罪与拐骗儿童罪的区分：前者有目的二，后者无目的二。

（3）本罪与绑架罪的区分：前者的目的二是出卖目的，后者的目的二是向第三人勒索财物的目的。

2.实行行为二：贩卖。

（1）贩卖行为（出卖行为，也即行为二）要成为本罪的实行行为，前提条件是，不存在行为一（绑架行为，实力控制行为）。结论：当贩卖行为是本罪的实行行为，那么，

% Source PDF physical page 240; see audit/source_coverage.json.
实施了贩卖行为，就成立本罪；卖掉，是既遂。

例1，甲捡了一个弃婴，然后想卖掉。由于捡拾行为不属于绑架行为，不属于拐卖罪的实行行为，此时只有贩卖行为才是实行行为。甲带着出卖目的，实施出卖行为，便成立拐卖儿童罪；如果卖掉了，便构成既遂。（2025年试题）

例2，甲以收养为目的，拐骗到一个儿童，构成拐骗儿童罪后，又觉得该儿童有点傻，养起来太累，便想卖掉。由于带着收养目的的拐骗行为不是拐卖儿童罪的实行行为，此时只有贩卖行为才是实行行为。甲带着出卖目的，实施出卖行为，便成立拐卖儿童罪；如果卖掉了，构成既遂。罪数上，根据司法解释，拐卖儿童罪吸收拐骗儿童罪，只定拐卖儿童罪。

例3，甲买到一个妇女做老婆，又嫌弃，想卖掉。带着收买目的的收买行为构成收买被拐卖的妇女罪，该收买行为不是拐卖妇女罪的实行行为，此时只有贩卖行为才是实行行为。甲带着出卖目的，实施出卖行为，便成立拐卖妇女罪；如果卖掉了，构成既遂。罪数上，根据法条规定，拐卖妇女罪吸收收买被拐卖的妇女罪，只定拐卖妇女罪。

（2）贩卖行为与赠与行为的区分：有无非法获利目的，为此，需要看收的钱的数额，能否评价为将妇女、儿童作为商品的对价。

例如，甲想卖自己的孩子。由于不存在绑架到手或拐到手的环节，此时只有贩卖行为才是实行行为。甲带着出卖目的，实施出卖行为，便成立拐卖儿童罪；如果卖掉了，构成既遂。如果甲收对方的钱只有2000元，则属于赠与（免费送），不属于贩卖，不构成拐卖儿童罪。

[注意1]司法解释规定：注意区分借送养之名出卖亲生子女与民间送养行为的界限。关键区分在于是否具有非法获利的目的。下列情形属于出卖亲生子女，应以拐卖儿童罪论处：第一，\CJKunderwave{将生育作为非法获利手段，生育后即出卖子女的}；第二，明知对方不具有抚养目的，或者根本不考虑对方是否具有抚养目的，为收取钱财将子女“送”给他人的；第三，\CJKunderwave{为收取明显不属于“营养费”“感谢费”的巨额钱财将子女“送”给他人的}。

[注意2］将亲生子女当作商品出卖，既构成拐卖儿童罪，也构成遗弃罪，想象竞合，择一重罪论处。（2016年第58题）

3.实行行为三：收买。

收买行为作为本罪的实行行为，必须带着出卖目的，也即是一种“进货”行为。收买被拐卖的妇女、儿童罪的收买行为只有收买的目的。结论：带着出卖目的，实施收买行为，便成立拐卖妇女、儿童罪；买到手，构成既遂。（2011年第54题）

4.实行行为四：拐骗。

这里的拐骗，主要是指欺骗，不包括绑架劫持行为。拐骗行为作为本罪的实行行为，必须带着出卖目的。拐骗儿童罪的拐骗行为没有出卖目的。结论：带着出卖目的，实施欺骗行为，便成立拐卖妇女、儿童罪；卖掉，构成既遂。

例1，电影《盲山》中，人贩子甲欺骗女大学生乙，以介绍工作为名，将乙带到收买者丙家，交给丙，自已偷偷溜掉。乙被丙控制。甲没有实力控制、绑架劫持乙，而是欺骗乙，仍构成拐卖妇女罪，并既遂。

例2，陈某（女）到偏僻山区旅游，夜宿王某家，王某想起其堂弟万某没有媳妇，于是第二天对陈某谎称带其坐公交车去观光，将陈某带至万某家。王某对万某讲，让陈某做万某的老婆，然后收取万某10块钱作为车费，坐公交车离去。王某属于拐骗行为，但是

% Source PDF physical page 241; see audit/source_coverage.json.
没有出卖目的，因为收取10块钱车费，不算将妇女作为商品出卖的对价，因此不构成拐卖妇女罪。王某没有拘禁陈某，也不构成非法拘禁罪。（2022年试题）

\InsertDiagram{figures/manual_table_p241_1.tex}

\subsubsection*{（二）法定刑升格条件}

1.结合犯：本罪+强奸罪=本罪（加重处罚）。对此不再数罪并罚。

[注意］法律没有规定这样的结合犯，也即本罪+强制猥亵罪=本罪（加重处罚）。因此，对这种情形应数罪并罚。（2010年第61题）

2.结合犯：本罪+引诱卖淫罪、强迫卖淫罪=本罪（加重处罚）。

3.结合犯：本罪+组织、运送他人偷越国（边）境罪=本罪（加重处罚）。

4.结果加重犯：“造成被拐卖的妇女、儿童或者其亲属重伤、死亡或者其他严重后果的”，加重处罚。该结果加重犯的要点与“强奸罪致人重伤、死亡”的要点相同。

（1）伤亡结果与拐卖行为之间具有直接因果关系。例如，甲拐卖妇女，妇女自杀，不属于这里的死亡结果，但可以“其他严重后果”论处。

（2）主观上是过失造成，或者为了实现拐卖目的故意造成。例如，妇女反抗时，为了制服妇女而对其实施暴力，致其重伤。如果不是出于拐卖目的，而是出于泄愤、报复等其他目的，故意将妇女打成重伤，则单独构成故意伤害罪，应与拐卖妇女罪并罚。（2010年第61题)

（3）重伤、死亡的被害人只包括被拐卖的妇女、儿童和前来阻挡的亲属，不包括其他前来阻挡的第三人。例如，甲拐卖妇女，行人乙前来阻挡，甲为了排除障碍，将乙打成重伤，进而顺利拐卖了妇女。甲构成拐卖妇女罪和故意伤害罪，并罚。

\subsection*{四、收买被拐卖的妇女、儿童罪}

收买被拐卖的妇女、儿童罪，是指不以出卖为目的，故意用财物收买被拐卖的妇女、儿童的行为。

1.既遂标准：买到手就既遂

例1，在接收到手前，拐卖者“一女二嫁”，将妇女又卖给了他人，行为人构成犯罪未遂。

例2（仙人跳），甲与妻子乙通谋，将乙“卖给”丙，获得钱财后，两人便逃离。甲乙构成诈骗罪，甲不构成拐卖妇女罪。\CJKunderwave{注意}：此时收买人构成对象不能犯，无罪，不应认定为收买被拐卖的妇女罪未遂。

2.罪数问题

（1）原则上，收买罪+后罪，应数罪并罚。常见后罪有拘禁、伤害、杀害、虐待、侮辱、猥亵、强奸、强迫卖淫、组织乞讨等。

（2）例外：收买罪+拐卖罪=拐卖罪。例如，甲收买了一位被拐卖的妇女，本想与之

% Source PDF physical page 242; see audit/source_coverage.json.
结婚，但又嫌弃，便卖掉。对甲只以拐卖妇女罪论处。如果甲收买时，就是想卖掉，然后卖掉，此时收买行为不构成收买被拐卖的妇女罪，整体就定一个正常的拐卖妇女罪。

（3）注意：收买罪+新罪+拐卖罪。结论：新罪如果能被拐卖罪吞并，则只定拐卖罪；如果不能被吞并，则应数罪并罚。

例如，甲带着结婚目的，收买了一位被拐卖的妇女，后实施了非法拘禁罪、强奸罪，然后又卖掉。由于非法拘禁罪、强奸罪能够被拐卖罪吞并，所以只定拐卖罪。如果甲的新罪是故意伤害罪（重伤），然后想卖掉并卖掉，则应定拐卖罪和故意伤害罪，数罪并罚。注意，此时不能定拐卖罪致人重伤，因为拐卖罪致人重伤是指拐的实行行为本身致人重伤。

3.处罚。收买被拐卖的妇女、儿童，对被买儿童没有虐待行为，不阻碍对其进行解救的，可以从轻处罚；按照被买妇女的意愿，不阻碍其返回原居住地的，可以从轻或者减轻处罚。

\subsection*{五、拐骗儿童罪}

第262条拐骗不满十四周岁的未成年人，脱离家庭或者监护人的，处五年以下有期徒刑或者拘役。

1.本罪侵害的法益是家长的监护权和儿童的人身自由、安全。只要侵害两种法益之一，就构成本罪。

2.行为方式：拐骗。这里的拐骗，既包括实力控制手段（绑架、劫持），也包括欺骗。

3.拐骗儿童罪与拐卖儿童罪的区分：有无出卖目的。判断方法：看有没有收买方。注意：两罪的关系是A与A+B的关系，B是出卖目的。两罪不是A与-A的关系。成立拐骗儿童罪，不要求特别目的。如果有出卖目的，则成立拐卖儿童罪。

[总结本节罪名]

1.罪名区分。本节罪名中，非法拘禁罪，绑架罪，拐卖妇女、儿童罪，拐骗儿童罪这四个罪名的客观行为都有实力控制行为，构成何罪，关键看带着什么目的实施实力控制。在此主要贯彻“行为与目的同时存在”原则，也即，行为的定性，由目的决定。

\InsertDiagram{figures/manual_table_p242_1.tex}

2.承继共犯问题。本节罪名中，非法拘禁罪，绑架罪，拐卖妇女、儿童罪，拐骗儿童罪这四个罪名的实行行为都有实力控制的方式。此时属于继续犯。既遂之后，有人参与进来，继续帮助实力控制的，构成承继的共同犯罪。但收买被拐卖的妇女、儿童罪不是继续犯。收买行为结束，该罪便结束。（2023年试题）

% Source PDF physical page 243; see audit/source_coverage.json.
例如，《亲爱的》电影里，黄渤角色的孩子被赵薇角色的丈夫（甲）拐骗到手，已经既遂，交给妻子乙（赵薇角色）养起来。若甲告知乙真相，则乙构成拐骗儿童罪的承继的共犯。

3.罪数问题

（1）原则上，实施两个行为，构成两个罪，应数罪并罚。

例1，收买罪+后罪，应数罪并罚。后罪诸如：包括非法拘禁、强奸等，均应并罚。

例2，拐骗儿童罪+后罪，应数罪并罚。后罪诸如：奸淫女童，构成强奸罪，并罚；向儿童父母勒索财物，构成绑架罪，并罚。

（2）拐卖罪+后罪，原则上数罪并罚。例外是结合犯：拐卖罪+强奸罪=拐卖罪（加重处罚)。

（3）三个拐卖类犯罪的内部关系：收买罪或拐骗罪+拐卖罪=拐卖罪。

\OriginalPanel{典型真题}

陈某（女）到偏僻山区旅游，夜宿王某家，王某想起其堂弟万某没有媳妇，于是第二天对陈某谎称带其坐公交车去观光，将陈某带至万某家。王某对万某讲，让陈某做万某的老婆，然后收取万某10块钱作为车费，坐公交车离去。万某扣留陈某，要求其当自己的老婆。陈某不同意，表示愿意给2万元钱，称“够你买个老婆了”。万某不同意，索要3万元，否则不放其离开。陈某被迫同意并给钱。万某释放陈某。下列说法正确的有？（2022年回忆版金题）\footnote{［答案］A项，王某将陈某骗至万某家，过程中没有现实剥夺陈某的人身自由，不属于非法拘禁行为。B项，王某没有收取对价，因此不构成拐卖妇女罪。因此，万某也不构成收买被拐卖的妇女罪。C项，陈某对万某称“够买老婆了”，不属于引起万某确切、具体的犯意，因此不构成教唆犯。D项，万某扣留陈某，属于拘禁行为，足以压制陈某的反抗，属于拘禁型抢劫，构成抢劫罪。本题答案为ABD。}

A.王某将陈某骗至万某家的行为，不属于非法拘禁行为

B.王某将陈某留至万某家的行为不构成拐卖妇女罪，万某扣押陈某的行为也不构成收买被拐卖的妇女罪

C.陈某给万某2万元，并称“够买老婆了”，属于教唆万某实施收买被拐卖的妇女罪的教唆犯

D.万某取得陈某的3万元，构成侵犯财产罪

\section{侵犯名誉权、人身权的犯罪}

[问题］甲女当众大骂第三者乙女是“狐狸精”“全家都是婊子”。甲女属于侮辱还是诽谤？（当众骂人案）

\subsection*{一、侮辱罪}

第246条以暴力或者其他方法公然侮辱他人或者捏造事实诽谤他人，情节严重的，处三年以下有期徒刑、拘役、管制或者剥夺政治权利。

前款罪，告诉的才处理，但是严重危害社会秩序和国家利益的除外。

% Source PDF physical page 244; see audit/source_coverage.json.
通过信息网络实施第一款规定的行为，被害人向人民法院告诉，但提供证据确有困难的，人民法院可以要求公安机关提供协助。[本条第3款为《刑法修正案（九）》所增设］

本罪的保护法益是他人的外部名誉。

1.侮辱对象：特定自然人

（1）必须特定具体。可以是一人，也可以是数人。泼妇当众骂街，不构成侮辱罪。

（2）不包括死人。但是通过侮辱死人侵害了死者家属名誉的，构成侮辱罪。

（3）不包括单位。但是通过侮辱单位侵害了特定自然人名誉的，构成侮辱罪。

2.侮辱行为

（1）暴力：这里的暴力是指程度较低的有形力，不包括程度较高的故意杀人、故意伤害行为。例如，当众打耳光、给对方头上泼粪、强迫他人从其跨下穿过等。

（2）其他方法：这是指非暴力的方法，包括使用语言、文字、图像等。例如，谩骂毁、张贴大字报、用电脑嫁接裸体图像并传播。

（3）侮辱行为必须公然进行。\CJKunderwave{公然性，是指能为多数人或不特定人所知晓}；但不要求被害人在现场。

（4）侮辱内容可以是真实事实，也可以是捏造的虚假事实。

\subsection*{二、诽谤罪}

诽谤罪，是指捏造并散布虚假事实，足以败坏他人名誉，情节严重的行为。本罪的保护法益是他人的外部名誉。

1.诽谤对象：特定自然人

（1）必须特定具体。可以是一人，也可以是数人。

（2）不包括死人。但是通过诽谤死人，侵害了死者家属名誉的，构成诽谤罪。

2.诽谤行为

（1）捏造虚假事实+公然散布。

（2）明知是捏造事实，故意公然散布。诽谤行为要求公然进行。\CJKunderwave{公然性，是指让多数}\CJKunderwave{人或不特定人知晓}。

\InsertDiagram{figures/manual_table_p244_1.tex}

［注意1」虚假事实不足以使人相信，不是诽谤，但有可能构成侮辱。例如，当众骂对方“脑子里都是粪便”“狐狸精”“全家人都是婊子生的”，不是诽谤，而是侮辱。所以，本节开头“当众骂人案”中，甲属于侮辱，而非诽谤。（2013年第16题）

[注意2]诽谤罪在成立要件上比侮辱罪更特殊。若同时触犯两罪，优先认定为诽谤罪。

% Source PDF physical page 245; see audit/source_coverage.json.
\subsection*{三、诬告陷害罪}

第243条捏造事实诬告陷害他人，意图使他人受刑事追究，情节严重的，处三年以下有期徒刑、拘役或者管制；造成严重后果的，处三年以上十年以下有期徒刑。

国家机关工作人员犯前款罪的，从重处罚。

不是有意诬陷，而是错告，或者检举失实的，不适用前两款的规定。

本罪的保护法益是人身权利，主要是人身自由，因为一旦被追究刑事责任，往往会坐牢。本罪的保护法益不是司法机关的司法秩序，因为本罪被规定在“侵犯人身权利、民主权利罪”中，没有规定在“妨害司法罪”中。\footnote{少数说认为，本罪的保护法益是司法机关的司法秩序。}

1.行为主体是已满16周岁的人。

国家机关工作人员犯本罪，从重处罚。本罪是不真正身份犯，国家机关工作人员是量刑身份。\CJKunderwave{注意}：报复陷害罪（第254条）的主体是国家机关工作人员，是真正身份犯。

2.行为对象是“他人”。

（1）“他人”必须是存在的并且特定的。

（2）“他人”只包括自然人，不包括单位。但是，诬告单位足以让司法机关怀疑确认某个自然人，成立本罪。（2017年第16题）

3.实行行为

（1）捏造事实，只是预备行为。实施诬告，才是实行行为，也即向公安、司法机关或有关国家机关告发。本罪的着手标准：向公安司法机关实施告发行为。

例如，甲写好诬告信，内容是狗蛋强奸了小芳。甲还没告发却把信弄丢了。乙捡到该信，明知内容是假的，却向公安机关告发。甲构成诬告陷害罪的犯罪预备。乙没有捏造事实这一环节，但有诬告陷害罪的实行行为。

（2）诬告内容：虚假的犯罪事实。

第一，甲有A罪事实，乙因为误解而告发有B罪事实，不是诬告。例如，甲有强制猥亵妇女行为，乙以为是强奸而告发强奸，不是诬告。

第二，甲有轻罪事实，乙告发，公安机关不立案，乙为了促使立案，告发甲犯了同性质的重罪，因为没有诬告陷害甲的故意，不构成诬告陷害罪。例如，乙告发甲犯了盗窃罪，公安不立案，乙便告发甲犯了抢劫罪。乙不构成诬告陷害罪。（2013年第59题）

第三，甲的犯罪事实存在，乙告发时，在犯罪情节上有程度差异，不是诬告。例如，乙告发甲贪污了100万元，经查，甲贪污了10万元。乙不属于诬告。

4.主观是故意，要求有诬告目的。不是有意诬陷，而是错告，或者检举失实的，不成立犯罪。

5.既遂标准：公安司法机关收到诬告材料，足以引起公安司法机关的调查活动。注意：不要求公安司法机关启动调查程序，更不要求实际追究了被害人的刑事责任。因为本罪的保护法益是他人的人身权（人身自由、名誉），当公安司法机关准备启动调查程序时，他人的名誉已经受到了侵害。（2006年第54题）

% Source PDF physical page 246; see audit/source_coverage.json.
\InsertDiagram{figures/manual_table_p246_1.tex}

\subsection*{四、刑讯逼供罪}

第247条（第一句）司法工作人员对犯罪嫌疑人、被告人实行\CJKunderwave{刑讯逼供}或者使用\CJKunderwave{暴力逼取证人证言}的，处三年以下有期徒刑或者拘役。（第二句）\CJKunderwave{致人伤残、死亡的，依}\CJKunderwave{照本法第二百三十四条（故意伤害罪）、第二百三十二条（故意杀人罪）的规定定罪从重}\CJKunderwave{处罚}。

\subsubsection*{（一）构成要件}

1.行为主体是司法工作人员。

（1）本罪是真正身份犯。例如，一般人私设公堂讯问他人，构成非法拘禁罪或故意伤害罪。超市保安或普通单位的保卫处工作人员，都不是司法工作人员。

（2）未受公安机关正式录用，受委托履行侦查职责的人员或者合同制民警，可以成为本罪主体。但是超市保安、农村治安联防员不是本罪主体。

2.刑讯的对象是犯罪嫌疑人、被告人。

（1）警察在预审时实施刑讯逼供，也成立本罪。

（2）犯罪嫌疑人、被告人实际上是否构成犯罪，对本罪的成立没有影响。

3.行为方式：刑讯+逼供。只有刑讯，没有逼供，或只有逼供，没有刑讯，均不构成刑讯逼供罪。刑讯是指肉刑或变相肉刑。没有实施肉刑或变相肉刑的诱供、逼供，不是刑讯逼供。例如，欺骗或威胁使用肉刑但没有使用的，不是刑讯逼供。

4.主观是故意，要求有逼供目的。至于为公（如尽快结案）还是为私（趁机挟私报复），在所不问。

\subsubsection*{（二）罪数问题}

1.第一句规定。例如，警察甲刑讯逼供犯罪嫌疑人乙，刑讯的实行行为是，故意重伤、杀害乙，也即为了刑讯逼供而使用暴力。甲整体一个行为，同时触犯刑讯逼供罪和故意伤害罪（重伤）、故意杀人罪，想象竞合，择一重罪论处，定故意伤害罪（重伤）或故意杀人罪。（2010年第58题）

2.第二句规定。这里的致人伤残、死亡，是指过失所为。也即，该规定将过失致人重伤、死亡的行为拟制为故意伤害罪、故意杀人罪。如果故意重伤、杀害，按照正常原则，就构成故意伤害罪、故意杀人罪。

\subsection*{五、暴力取证罪}

本罪条文是第247条，与刑讯逼供罪同处一个条文。

1.行为主体是司法工作人员。本罪是真正身份犯。

2.取证对象是证人。这里的“证人”是广义的证人，既包括被害人、鉴定人、民事诉讼中的证人，也包括不具有作证资格的人、不知道案件真相的人，还包括年幼者、精神病患者。

3.暴力取证。使用胁迫、欺骗、测谎仪等方式，不属于暴力取证。

% Source PDF physical page 247; see audit/source_coverage.json.
\subsection*{六、虐待被监管人罪}

第248条监狱、拘留所、看守所等监管机构的监管人员对被监管人进行殴打或者\CJKunderwave{体}\CJKunderwave{罚虐待}，情节严重的，处三年以下有期徒刑或者拘役；情节特别严重的，处三年以上十年以下有期徒刑。\CJKunderwave{致人伤残、死亡的}，依照本法第二百三十四条、第二百三十二条的规定定罪从重处罚。

\CJKunderwave{监管人员指使被监管人殴打或者体罚虐待其他被监管人}的，依照前款的规定处罚。

1.被监管人，包括：第一，刑事未决犯和已决犯；第二，被行政拘留的人；第三，被强制戒毒的人。

2.“体罚虐待”，是指体罚或者虐待，而非仅限于体罚式虐待（2019年试题）。对被监管人实施性虐待，触犯强制猥亵罪的，属于想象竞合，择一重罪论处。

3.法律拟制。本罪第1款的最后一句“致人伤残、死亡的，定故意伤害罪、故意杀人罪”。这里的致人伤残、死亡，是指过失所为。也即，该规定将过失致人重伤、死亡的行为拟制为故意伤害罪、故意杀人罪。如果故意重伤、杀害，按照正常原则，就构成故意伤害罪、故意杀人罪。（2012年第16题）

\subsection*{七、非法侵入住宅罪}

1.保护法益。多数说认为，本罪的保护法益是居住者居住的安宁状态。并且，本罪是实害犯，不是危险犯。成立犯罪，要求造成实害结果，也即现实地侵害了居住者的居住安宁状态。少数说认为，本罪的保护法益是住宅的进入许可权。（2022年主观题）

2.住宅：供人起居生活的场所。

3.行为方式。（1）作为方式：身体非法侵入。（2）不作为方式：合法进入，拒不退出。

\subsection*{八、侵犯公民个人信息罪}

第253条之一\CJKunderwave{违反国家有关规定，向他人出售或者提供公民个人信息}，情节严重的，处三年以下有期徒刑或者拘役，并处或者单处罚金；情节特别严重的，处三年以上七年以下有期徒刑，并处罚金。

\CJKunderwave{违反国家有关规定，将在履行职责或者提供服务过程中获得的公民个人信息，出售或}\CJKunderwave{者提供给他人的}，依照前款的规定从重处罚。

窃取或者以其他方法非法获取公民个人信息的，依照第一款的规定处罚。

单位犯前三款罪的，对单位判处罚金，并对其直接负责的主管人员和其他直接责任人员，依照各该款的规定处罚。[本条经《刑法修正案（九）》修正]

司法解释要点\footnote{2017年6月1日最高人民法院、最高人民检察院《关于办理侵犯公民个人信息刑事案件适用法律若干问题的解释》；2023年9月20日最高人民法院、最高人民检察院、公安部《关于依法惩治网络暴力违法犯罪的指导意见》。}：

1.本罪的“公民个人信息”，是指\CJKunderwave{以电子或者其他方式记录}的能够单独或者与其他信息结合识别特定自然人身份或者反映特定自然人活动情况的各种信息，包括姓名、身份证件号码、通信通讯联系方式、住址、账号密码、财产状况、行踪轨迹等。例如，甲男整天用望远镜偷窥对面楼层的女邻居洗澡。甲男不构成侵犯公民个人信息罪。

% Source PDF physical page 248; see audit/source_coverage.json.
2.违反国家有关规定，通过购买、收受、交换等方式获取公民个人信息，或者在履行职责、提供服务过程中收集公民个人信息的，属于本罪的“以其他方法非法获取公民个人信息”。

3.组织“人肉搜索”，违法收集并向不特定多数人发布公民个人信息，情节严重，构成侵犯公民个人信息罪。

\subsection*{九、强迫劳动罪}

第244条以暴力、威胁或者限制人身自由的方法强迫他人劳动的，处三年以下有期徒刑或者拘役，并处罚金；情节严重的，处三年以上十年以下有期徒刑，并处罚金。

明知他人实施前款行为，为其招募、运送人员或者有其他协助强迫他人劳动行为的，依照前款的规定处罚。

单位犯前两款罪的，对单位判处罚金，并对其直接负责的主管人员和其他直接责任人员，依照第一款的规定处罚。

1.行为主体和对象。本罪主体包括自然人和单位。根据《刑法修正案（八）》的规定，本罪的对象由“职工”改为“他人”。

2.行为类型。本罪的行为类型有两种：

（1）直接强迫劳动，包括暴力、威胁、限制人身自由的方法。非法拘禁罪是剥夺人身自由，强迫劳动罪是限制人身自由。例如，不许私自外出。如果采取剥夺人身自由的方式强迫劳动，就成立非法拘禁罪和强迫劳动罪的想象竞合犯。

（2）协助强迫劳动。\CJKunderwave{注意}：这种帮助行为没有被正犯化。也即，在定罪条件上，仍需遵守共犯从属性。也即，如果实行者无罪，这种协助行为也无罪。例如，甲为乙招募、运送了3名人员，但乙没有强迫他们劳动。甲不构成犯罪。

\subsection*{十、雇用童工从事危重劳动罪（第 244 条之一）}

本罪是指，违反劳动法，雇用未满16周岁的人从事危重劳动，情节严重。

1.危重劳动包括：超强度体力劳动，在危险环境劳动，如高空、井下、易燃易爆环境。

2.雇用，包括不支付报酬的情形。这是因为，支付报酬，构成本罪；不支付报酬，更构成本罪。

3.故意。例如，甲误以为刑法规定，本罪的童工是指十四岁以下的儿童，便雇用15周岁的人从事危重劳动。甲对事实没有认识错误，存在故意。甲有法律认识错误，但有违法性认识可能性，所以甲构成本罪。（2025年试题）

\section{侵犯婚姻家庭的犯罪}

[问题］甲女明知乙男的妻子是军人，仍和乙男同居。甲女构成何罪？（同居案）

\subsection*{一、重婚罪}

第258条有配偶而重婚的，或者明知他人有配偶而与之结婚的，处二年以下有期徒刑或者拘役。

% Source PDF physical page 249; see audit/source_coverage.json.
本罪的保护法益是一夫一妻的婚姻制度和合法婚姻关系。

1.行为主体与行为类型。

（1）已婚者：同时维持两份婚姻。

这两份婚姻，包括：第一，法律婚姻+法律婚姻；第二，法律婚姻+事实婚姻。\footnote{事实婚姻，是指二人公开同居，并且以夫妻名义对外相称，周围群众公认二人是夫妻。在此强调事实婚姻，不意味着承认事实婚姻是有效婚姻，而是因为它破坏了法律婚姻，才承认它的存在。民法关注效力问题，刑法关注事实问题。}

这两份婚姻，不包括：第一，事实婚姻+事实婚姻；\footnote{这是因为，法条规定“有配偶而重婚”，这里的“有配偶”是指有合法配偶。}第二，法律婚姻+同居。\footnote{提示：同居，是指共同生活。法律婚姻是指登记结婚。}

(2）相婚者（事前单身者）：明知对方有配偶而与之结婚（法律婚姻或事实婚姻）。这种人实际上只有一份婚姻，但因为破坏他人的合法婚姻，因此构成重婚罪。（2012年第55题）

2.性质。重婚罪不仅侵犯了合法配偶的婚姻关系（个人法益），还侵害了一夫一妻的婚姻制度（社会法益）。因此，即使合法妻子同意丈夫重婚，丈夫也构成重婚罪。也因此，重婚罪不是告诉才处理的案件，而是公诉案件。

\subsection*{二、破坏军婚罪}

第259条明知是现役军人的配偶而与之同居或者结婚的，处三年以下有期徒刑或者拘役。

利用职权、从属关系，以胁迫手段奸淫现役军人的妻子的，依照本法第二百三十六条（强奸罪）的规定定罪处罚。

1.行为类型：行为人与军人的配偶结婚或同居。

例如，甲男（军人）+乙女（非军人），是夫妻。

模型一：A男+乙女，结婚。

结论：（1）A男构成重婚罪，因为明知乙女有配偶仍与乙女结婚。A男还构成破坏军婚罪，因为明知乙女有配偶且配偶是军人，仍与乙女结婚。两罪此时是法条竞合关系，最终定破坏军婚罪。（2）乙女构成重婚罪。

模型二：甲男+B女，结婚。

结论：（1）B女构成重婚罪，不构成破坏军婚罪，因为是和军人结婚，不是和军人的配偶结婚。（2）甲男构成重婚罪。

2.行为方式：结婚或同居。这里的结婚包括法律婚姻和事实婚姻。这里的同居，包括公开同居和秘密同居。例如，本节开头“同居案”中，甲女构成破坏军婚罪，乙男不构成破坏军婚罪，也不构成重婚罪，因为重婚罪中不包括同居。

注意：行为方式不包括通奸。同居和通奸的区别在于：前者具有共同生活的特征，而后者没有。如果将这里的同居解释为包括通奸，属于类推解释，应当禁止。

3.第2款：“利用职权、从属关系，以胁迫手段奸淫现役军人的妻子的，以强奸罪论处”。该款是注意规定。也即只有行为符合强奸罪的构成要件时，才能以强奸罪论处。因此，该款中的“胁迫手段”是指强奸罪中的胁迫手段。

% Source PDF physical page 250; see audit/source_coverage.json.
\subsection*{三、暴力干涉婚姻自由罪}

第257条以暴力干涉他人婚姻自由的，处二年以下有期徒刑或者拘役。

犯前款罪，\CJKunderwave{致使被害人死亡的}，处二年以上七年以下有期徒刑。

第一款罪，告诉的才处理。

本罪侵害的法益是他人的婚姻自由权。

1.行为主体与行为对象：不要求有特定关系。例如，甲暴力干涉自己女儿的婚姻自由，成立本罪；甲暴力干涉邻居女儿的婚姻自由，也成立本罪。

2.暴力：这里的暴力是指程度较低的暴力。例如，捆绑、拘禁，但是不包括极为轻微的暴力（如打一耳光）。当然，如果暴力程度达到重伤，则与故意伤害罪想象竞合，择一重罪论处，按故意伤害罪论处。

3.婚姻自由。（1）既包括结婚自由，也包括离婚自由。（2）不包括恋爱自由和分手自由。（3）既包括强制甲与乙结婚或离婚，也包括强制甲与自己结婚。注意：由于夫妻关系比较特殊，丈夫强制妻子跟自己离婚或不同意跟自己离婚，使用暴力，不成立本罪，可成立虐待罪。

4.“犯前款罪，致使被害人死亡的”，这是该罪的结果加重犯。这里的“死亡”包括被害人自杀。（2006年第56题）

\subsection*{四、虐待方面的犯罪}

1.虐待罪

第260条虐待家庭成员，情节恶劣的，处二年以下有期徒刑、拘役或者管制。

犯前款罪，致使被害人重伤、死亡的，处二年以上七年以下有期徒刑。

第一款罪，告诉的才处理，但被害人没有能力告诉，或者因受到强制、威吓无法告诉的除外。[本条第3款经《刑法修正案（九）》修正]

（1）本罪行为主体和行为对象是家庭成员。该“家庭成员”可以作扩大解释，可包括长年共同生活的管家、保姆和有事实婚姻关系的“夫妻”。

（2）第260条第2款中的“致使被害人死亡”：第一，包括被害人自杀。虽然被害人自杀与虐待行为没有直接因果关系，仍视为结果加重犯（2006年第56题）。第二，只包括被害人自杀，不包括近亲属自杀。（2026年试题）

（3）在构成虐待罪的情况下，又故意杀人（如故意饿死），故意伤害导致重伤的，则构成故意杀人罪、故意伤害罪（重伤），与虐待罪并罚。（2025年试题）

（4）采取禁闭方式虐待，同时构成虐待罪和非法拘禁罪，想象竞合，择一重罪论处。（2025年试题）

2.虐待被监护、看护人罪

第260条之一对未成年人、老年人、患病的人、残疾人等\CJKunderwave{负有监护、看护职责的人}\CJKunderwave{虐待被监护、看护的人}，情节恶劣的，处三年以下有期徒刑或者拘役。

单位犯前款罪的，对单位判处罚金，并对其直接负责的主管人员和其他直接责任人员，依照前款的规定处罚。

有第一款行为，同时构成其他犯罪的，依照处罚较重的规定定罪处罚。[本条为《刑法修正案（九）》所增设]

% Source PDF physical page 251; see audit/source_coverage.json.
\InsertDiagram{figures/manual_table_p251_1.tex}

[总结]属于自诉案件的五个罪名

\InsertDiagram{figures/manual_table_p251_2.tex}
